%
% 1.1 what is a language model
% 1.2 introduction of transformers
% 1.3 large language models
% 2.1 transformers based models for code analysis, such as codebert, codet5
% 2.2 large language model for code analysis and HPC analysis
% 1.4 what is missing for LM or LLM for HPC
%\todo{A figure about large language model may be needed to help explain its components}


Language models (LMs) are machine learning models designed to comprehend and generate human language. They can be used to facilitate natural and intuitive interactions between humans and machines. 
Early generations of LMs, using recurrent neural networks (RNNs),  showed inspiring results for various natural language processing (NLP) tasks.  
A transformative evolution by the Transformer \cite{vaswani2017attention} reveals remarkable potentials of LMs.
Introduced by Vaswani \textit{et al.}, transformer
models utilize the attention mechanism to capture the dependencies between all words in an input sentence, irrespective of their positions. Compared to RNNs, transformers process data in parallel rather than sequentially and significantly improve the efficiency of model training and inference. %Transformer architecture has become a cornerstone in the field of NLP. 
Transformers further enables the inauguration of the large language models (LLMs).
%Large language models (LLMs) are a type of language models that are trained on a vast amount of data. 
Compared to LMs, LLMs are trained on a vast amount of data and possess parameter counts on the order of billions or more, allowing them to generate more detailed and nuanced responses. Examples of LLMs include OpenAI's GPT-3, GPT-4 and Google's BARD. Nowadays, LLMs have shown remarkable capabilities in NLP tasks like translation, question answering, and text generation.

%The evolution of language models has been transformative, migrating from recurrent neural networks (RNNs) to the more advanced transformer models~\cite{vaswani2017attention}, with remarkable success across various natural language processing (NLP) tasks. 
% However, this venture comes with its unique set of challenges:


% \begin{table}[htbp]
% \centering
% \caption{Language models, associated training data and tasks}
% \label{tab-modelExample}
% \begin{threeparttable}
% \begin{tabular}{|c|c|c|c|c|c|}
% \hline
% \multicolumn{2}{|c}{\textbf{Model}} &
% \multicolumn{2}{|c|}{\textbf{Training Data}} &
% \multirow[t]{2}{*}{\textbf{TokenLimit}} &
% \multirow[t]{2}{*}{\textbf{Availability}} \\
% \cline{1-4}
%  \textbf{Name} & \textbf{Size} & \textbf{Type} & \textbf{Size} &  & \\ \hline
%  BERT & 340M & Text & 3.5B words & 512 & code, weights\\ \hline
%  CodeBERT & 125M & Mixed & \makecell{2.1M (bimodal)\\6.4M (unimodal)} & 512 & code, weights\\ \hline
%  GraphCodeBERT & 110M & Code & 2.3M functions  & 512 & code, weights\\ \hline
%  CodeT5 & 770M & Code &  8.35M instances  & 512 & code, weights\\ \hline
%  GPT-3 & 175B & Mixed &  500B words  & 4096 & API\\ \hline
% Cerebras-GPT & 0.11B$\sim$13B & Mixed &  800GB  & 2048 & code, weights\\ \hline
% LLaMA & 7$\sim$65B & Mixed &  1.4T tokens  & 4096 & code\\ \hline
% StarCoder & 15B & Code &  1T tokens  & 8192  & code, weights\\ \hline
% Dolly 2.0 & 3$\sim$12B & Text & 15k inst./response  & 2048 & code, weights\\ \hline
% BARD & 1.6B & Mixed & 1.56T words & 1000 & API\\ \hline
% GPT-4 & $>$ 1T & Mixed & undisclosed & 8k/32k & API\\ \hline
% \end{tabular}
% %token limits for gpt-4 and gpt-4-32k are 8192 and 32768 respectively.


% %\begin{tablenotes}\footnotesize
% %\item[*]NLP: natural language processing, PLP: programming language processing \newline
% %LM: language model, LLM: large language model \newline
% %NLT: natural language text, PLT: programming language text \newline

% %\end{tablenotes}
% \end{threeparttable}

% \end{table}



% Please add the following required packages to your document preamble:
% \usepackage{multirow}
% Please add the following required packages to your document preamble:
% \usepackage{multirow}
% Please add the following required packages to your document preamble:
% \usepackage{multirow}

%table in a google sheet: 
%%%  https://docs.google.com/spreadsheets/d/1oS3aB7QsL7dkcCm6Lkv_IviBrFzlKiw_3VM0V8LaNOM/edit?usp=sharing
% Please add the following required packages to your document preamble:
% \usepackage{multirow}
%\todo{\textcolor{blue}{added in table 1} add  NVIDIA’s Megatron: }
\vspace{-12pt}
\begin{table}[]
\caption{Language models, associated training data and tasks}
\label{tab-modelExample}
\begin{tabular}{|ccc|cc|c|c|}
\hline
\multicolumn{3}{|c|}{\textbf{Model}} &
  \multicolumn{2}{c|}{\textbf{Training data}} &
  \multirow{2}{*}{\textbf{\begin{tabular}[c]{@{}c@{}}Token\\  Limit\end{tabular}}} &
  \multirow{2}{*}{\textbf{Avail.}} \\ \cline{1-5}
\multicolumn{1}{|c|}{\textbf{Name}} &
  \multicolumn{1}{c|}{\textbf{Released}} &
  \textbf{Size} &
  \multicolumn{1}{c|}{\textbf{Type}} &
  \textbf{Size} &
   &
   \\ \hline
\multicolumn{1}{|c|}{BERT}                     & \multicolumn{1}{c|}{2018/10}       & 340M      & \multicolumn{1}{c|}{Text}  & 3.5B words      & 512    & Weights \\ \hline
\multicolumn{1}{|c|}{CodeBERT} &
  \multicolumn{1}{c|}{2020/11} &
  125M &
  \multicolumn{1}{c|}{Mixed} &
  \begin{tabular}[c]{@{}c@{}}2.1M(bimodal)\\  6.4M (unimodal)\end{tabular} &
  512 &
  Weights \\ \hline
  \multicolumn{1}{|c|}{Megatron}            & \multicolumn{1}{c|}{2021/04}       & 1T      & \multicolumn{1}{c|}{Text}  & 174 GB  & 512    & Weights \\ \hline
\multicolumn{1}{|c|}{GraphCodeBERT}            & \multicolumn{1}{c|}{2021/05}       & 110M      & \multicolumn{1}{c|}{Code}  & 2.3M functions  & 512    & Weights \\ \hline
\multicolumn{1}{|c|}{CodeT5}                   & \multicolumn{1}{c|}{2021/11}       & 770M      & \multicolumn{1}{c|}{Code}  & 8.35M instances & 512    & Weights \\ \hline
\multicolumn{1}{|c|}{GPT-3}                    & \multicolumn{1}{c|}{2022/03/15} & 175B      & \multicolumn{1}{c|}{Mixed} & 500B tokens           & 4096   & Weights \\ \hline
\multicolumn{1}{|c|}{LLaMA}                    & \multicolumn{1}{c|}{2023/02/24}   & 7$\sim$65B     & \multicolumn{1}{c|}{Mixed} & 1.4T tokens      & 4096   & Weights*          \\ \hline
\multicolumn{1}{|c|}{GPT-4}                    & \multicolumn{1}{c|}{2023/03/14}   & 1T        & \multicolumn{1}{c|}{Mixed} & undisclosed     & 8k/32k & API*           \\ \hline
\multicolumn{1}{|c|}{BARD}                     & \multicolumn{1}{c|}{2023/03/21}   & 1.6B      & \multicolumn{1}{c|}{Mixed} & 1.56T words            & 1000   & API           \\ \hline
\multicolumn{1}{|c|}{Cerebras-GPT}             & \multicolumn{1}{c|}{2023/03/28}    & 0.11$\sim$13B & \multicolumn{1}{c|}{Text} & 800 GB           & 2048   & Weights \\ \hline
\multicolumn{1}{|c|}{Dolly 2.0}                & \multicolumn{1}{c|}{2023/04/12}   & 3$\sim$12B     & \multicolumn{1}{c|}{Text}  & 15k instr./resp.             & 2048   & Weights \\ \hline
\multicolumn{1}{|c|}{\begin{tabular}[c]{@{}c@{}}StarCoder\end{tabular}} & \multicolumn{1}{c|}{2023/05/4}   & 15B       & \multicolumn{1}{c|}{Code}  & 1T tokens            & 8192   & Weights \\ \hline
\multicolumn{1}{|c|}{\begin{tabular}[c]{@{}c@{}} StarChat-Alpha\end{tabular}} & \multicolumn{1}{c|}{2023/05/4}   & 16B       & \multicolumn{1}{c|}{Code}  & 31k instr./resp.            & 8192   & Weights \\ \hline
\end{tabular}
\end{table}
\vspace{-15pt}

%\subsection{Transformers and LLMs}
%\subsection{LM/LLM Tasks}
% \vspace{-10pt}
Table~\ref{tab-modelExample} shows some example language models and their release dates, sizes, training data, input token length limits, and availability. LLaMA~\cite{touvron2023llama}'s weights can be obtained after filling out some form. GPT-4 has a waiting list to use its API. At the time of writing this paper, we have not yet obtained its access.
 
%Deploying LM applications could be complex as it requires an understanding of datasets, machine learning frameworks, models, and other technical aspects. This complexity is further amplified by the rapid evolution and expansion of LLMs. Hugging Face has become the leading platform for the machine-learning community to share and explore cutting-edge models and datasets. It supports all the mainstream machine-learning frameworks and offers thoughtfully designed APIs for the effective development of LM applications.  Hugging Face simplifies the process of deploying sophisticated LMs, making state-of-the-art NLP technology more accessible and manageable for developers and researchers.

% \textit{at least one billion or more parameters.}}
%\subsection{LM and LLM for Code}
LMs are trained mainly by text data with a focus on NLP.  The sources of the training data mainly come from books, web content, newspapers, scientific articles, and  other text data in various natural languages.    
 Latest LLMs have demonstrated rich skill sets in NLP including text prediction, common sense reasoning, reading comprehension, translation and question answering.

 
There has been a keen interest in deploying NLP techniques to programming language processing (PLP) tasks, such as code summarization, code generation, and code similarity analysis \cite{chen2022multi,flynn2022finding}. 
Previous studies have demonstrated successful applications of traditional language models to PLP tasks, showing the feasibility of this approach~\cite{devlin2018bert}.
%Previous research has extended transformer-based language models to the PLP domain. 
CodeBERT \cite{feng2020codebert}, for example, is a transformer-based model trained with a diverse range of programming languages and can be used for a variety of programming-related tasks. Similarly, CodeT5 \cite{wang2021codet5} is a variant of Google's T5 language model, trained specifically on code datasets to perform advanced programming tasks like code completion, bug detection, and code summarization.
Lately, StarCoder \cite{li2023starcoder}, a 15B parameter model trained with 1 trillion tokens sourced from a large collection of permissively licensed GitHub repositories, is developed to be a Large Language Model mainly for code generation or completion. 
StarChat-Alpha is a GPT-like chat model fine-tuned from StarCoder to act as a helpful coding assistant. 

% \color{red}{add starcoder}
% \textcolor{blue}{add starcoder; add HPCFAIR}
\vspace{-10pt}
\subsection{LMs for HPC}
With the recent breakthroughs in Generative Pretrained Transformer (GPT) large language models \cite{brown2020language}, it has become increasingly intriguing to explore the application of large language models (LLMs) for HPC tasks.
%Despite the significant advances made by Large Language Models (LLMs) in comprehending and generating both natural and programming languages, 
However, their deployment in the HPC domain is still relatively unexplored. This venture comes with various challenges, including:
\vspace{-4pt}

%\todo{\textcolor{blue}{added example} although there are a lot HPC code in ..., but not reprocessed ...}
\begin{enumerate}
    \item Pipelines: Traditional language model frameworks like Hugging Face were designed to support natural language processing or compute vision problems. Expanding LMs to any new domain, including HPC,  requires the addition of new pipelines designed to finish domain-specific tasks.   
    \item Datasets: The HPC domain encompasses an extensive amount of code spanning various fields, including biology and climate modeling. However, preparing this data for machine learning training, such as labeling parallelizable loops in HPC programs for parallelism detection, presents significant challenges. The scarcity of ready-to-use, pre-labeled HPC datasets poses a particular obstacle for training language models, especially large ones, highlighting the need for more shared resources in the community.
    \item Pre-processing: Pre-processing in the context of LMs for HPC typically involves the conversion of source files into a sequence of tokens. However, the direct application of NLP tokenizers to code can be sub-optimal. For instance, an NLP tokenizer might split a variable name into two tokens, a scenario that is not desirable for PLP analysis.  Also, models designed for processing source code may use graph representations, such as abstract syntax trees, to have better performance.
    \item Input size limit: Language models often have limited input token lengths (such as 512 to a few thousand of tokens). HPC tasks often involve processing large-scale software packages with millions of lines of source code. 
%    \item Fine-tuning:  To fully harness the capabilities of state-of-the-art LLMs, fine-tuning these models with HPC-specific data is necessary. This process demands a thorough understanding and expertise in handling LLMs. 
    \item Evaluation: There is a pressing need for standardized and reproducible evaluation of different models in the context of various HPC tasks, using metrics suitable for domain-specific requirements. 
    %This is particularly challenging given the stochastic nature of language models. The resulting leaderboards can facilitate communication regarding both the state-of-the-art and areas that are lacking. 
\end{enumerate}
% \vspace{-10pt}

